% ============================================================================
% Fichier  : partie6/P06_C18_rapport.tex : FORN (LPI342)
% Sources : ISO/IEC 27042 ; ACPO. Opération TSANAGA.
% ============================================================================

\chapter{Le Rapport Forensique}
\label{chap:rapport}
\chaptermark{Le rapport forensique}

\epigraph{%
    « Le rapport est la seule partie de l'enquête que le juge verra. »%
}{Principe du cours}

\niveauChapitre{Intermédiaire}{Toute l'enquête~; Chapitre~\ref{chap:reconstitution}.}

\begin{boxobjectifs}
\begin{itemize}
    \item Structurer un rapport forensique clair et complet.
    \item Séparer méthode, constatation et interprétation.
    \item Écrire pour un lecteur non technique sans perdre la rigueur.
    \item Documenter la reproductibilité.
    \item Assumer les limites et les NEK.
\end{itemize}
\end{boxobjectifs}

\bigskip

% ============================================================================
\section*{Question directrice : comment transformer des mois d'enquête en un document que le juge comprendra et pourra opposer ?}
% ============================================================================

L'enquête ne vaut que par son rapport. Il doit être compris par un
magistrat non technicien, et résister à la contestation d'un expert adverse.
Ce chapitre le construit.

\medskip
\textbf{Plan du chapitre.}
\begin{enumerate}
    \item La structure (section~\ref{sec:c18-structure}).
    \item Méthode, fait, interprétation (section~\ref{sec:c18-separer}).
    \item Écrire pour être compris (section~\ref{sec:c18-ecrire}).
    \item Reproductibilité et limites (section~\ref{sec:c18-limites}).
\end{enumerate}

\separateur

% ============================================================================
\section{La Structure}
\label{sec:c18-structure}
% ============================================================================

\begin{tableaucours}{L{4.0cm}L{8.2cm}}{%
    \tableauHeader{Partie} &
    \tableauHeader{Contenu}
}
Synthèse & Ce qui est établi, en termes simples, pour le magistrat \\
Cadre et mission & Réquisition, périmètre, pièces, dates \\
Méthode & Normes suivies, outils, actes, condensats \\
Constatations & Les faits établis, par domaine \\
Interprétation & Analyse, niveaux de confiance, hypothèses \\
Limites et NEK & Ce qui n'a pu être établi \\
Annexes & Chaîne de possession, condensats, journaux \\
\end{tableaucours}
\vspace{-0.8em}
\begin{center}{\small\itshape Tableau~18.1. Structure d'un rapport forensique}\end{center}

% ============================================================================
\section{Méthode, Fait, Interprétation}
\label{sec:c18-separer}
% ============================================================================

\begin{boxalert}
\textbf{Violation typique.} Mêler le fait et l'opinion (« le fichier a été
copié, donc NDONGO est un voleur »).

\smallskip
\textbf{Règle~:} on sépare nettement la \emph{méthode} (comment on a procédé),
la \emph{constatation} (le fait établi) et l'\emph{interprétation} (ce qu'on
en déduit, avec un niveau de confiance). Le juge doit pouvoir distinguer les
trois d'un coup d'œil.
\end{boxalert}

% ============================================================================
\section{Écrire pour être Compris}
\label{sec:c18-ecrire}
% ============================================================================

Le rapport s'adresse d'abord à un non-technicien. On explique les termes, on
illustre, on va du simple au détaillé~: une synthèse accessible en tête, le
détail technique ensuite, les preuves en annexe. La clarté n'est pas l'ennemie
de la rigueur~: elle la sert.

% ============================================================================
\section{Reproductibilité et Limites}
\label{sec:c18-limites}
% ============================================================================

\begin{boxscenario}{Opération TSANAGA. Le rapport remis}
\textbf{Constatation.} Le rapport de M.~BIKORO établit la copie (disque),
l'envoi (réseau et nuage), la négociation (mobile), et déclare le second
conteneur comme NEK. Chaque constat cite sa méthode et son condensat~; la
chaîne de possession est en annexe.

\medskip
\textbf{Analyse} (confiance élevée). Un expert adverse pourrait reproduire
l'analyse et obtenir le même résultat~: c'est la marque d'un rapport
opposable. Les limites assumées (le second conteneur) renforcent sa
crédibilité.
\end{boxscenario}

\begin{boxmetier}{Expert judiciaire}
\textbf{Écrire pour le contradictoire.} L'expert rédige en imaginant un
contradicteur~: chaque affirmation doit pouvoir être vérifiée, chaque méthode
reproduite. Un rapport qui anticipe la contestation la désamorce.
\end{boxmetier}

\separateur

\begin{boxresume}
\textbf{Ce qu'il faut retenir de ce chapitre~:}
\begin{itemize}
    \item \textbf{Le rapport} est la seule partie que le juge verra.
    \item \textbf{On sépare} méthode, constatation, interprétation.
    \item \textbf{On écrit} pour un non-technicien, sans perdre la rigueur.
    \item \textbf{On documente} la reproductibilité.
    \item \textbf{On assume} les limites et les NEK.
\end{itemize}
\end{boxresume}

\begin{boxinfo}
\textbf{Points de contrôle de l'élève.} Je sais~: structurer un rapport~;
séparer méthode, fait, interprétation~; écrire clair~; documenter la
reproductibilité~; assumer les limites.
\end{boxinfo}

\bigskip

\begin{boxexo}[title={\ding{45}\quad Exercice~18.1. Structure \niveaubase}]
Citez les parties d'un rapport forensique.
\end{boxexo}

\begin{boxexo}[title={\ding{45}\quad Exercice~18.2. Trois registres \niveaubase}]
Distinguez méthode, constatation, interprétation.
\end{boxexo}

\begin{boxexo}[title={\ding{45}\quad Exercice~18.3. Lecteur \niveaubase}]
À qui s'adresse d'abord le rapport~?
\end{boxexo}

\begin{boxexo}[title={\ding{45}\quad Exercice~18.4. Mêler \niveaumoyen}]
Pourquoi ne pas mêler le fait et l'opinion~?
\end{boxexo}

\begin{boxexo}[title={\ding{45}\quad Exercice~18.5. Reproductible \niveaumoyen}]
Comment rendre un rapport reproductible~?
\end{boxexo}

\begin{boxexo}[title={\ding{45}\quad Exercice~18.6. NEK \niveaumoyen}]
Pourquoi déclarer le second conteneur comme NEK dans le rapport~?
\end{boxexo}

\begin{boxexo}[title={\ding{45}\quad Exercice~18.7. Synthèse \niveaumoyen}]
Rédigez la synthèse (pour le magistrat) de l'affaire TSANAGA.
\end{boxexo}

\begin{boxexo}[title={\ding{45}\quad Exercice~18.8. Contradictoire \niveaumoyen}]
Comment écrire en anticipant la contestation~?
\end{boxexo}

\begin{boxexo}[title={\ding{45}\quad Exercice~18.9. Rapport complet \niveauexpert}]
Rédigez le plan détaillé du rapport TSANAGA, partie par partie.
\end{boxexo}

\begin{boxexo}[title={\ding{45}\quad Exercice~18.10. Constatation \niveauexpert}]
Rédigez une constatation complète (copie du fichier), méthode et condensat
inclus.
\end{boxexo}

\bigskip

\begin{boxinfo}
\textbf{Pour aller plus loin.}
\begin{itemize}
    \item ISO/IEC 27042~\cite{iso27042}~; ACPO~\cite{acpo2012}.
\end{itemize}
\end{boxinfo}

\begin{boxinfo}
\textbf{Transition.} Le rapport se défend à l'oral~: le
Chapitre~\ref{chap:tribunal} traite de la présentation devant un tribunal.
\end{boxinfo}
